\documentclass[reqno,11pt]{amsart}
\usepackage{graphicx}
\usepackage{amsfonts}
\usepackage{amssymb}
\usepackage{amsmath}
\usepackage{amsthm}
\usepackage{latexsym}
\usepackage{enumerate}
\usepackage{fullpage}
\usepackage[usenames,dvipsnames]{color}

\pagestyle{empty}
\thispagestyle{empty}

\begin{document}
\begin{flushright} (Robert Kohn) \end{flushright}
\begin{flushright} (3/27/26)\end{flushright}
\vspace{0.2in}
\begin{center}
{\bf Portfolio Problems for Test 2}
\vspace{0.2in}
\end{center}

\noindent\textbf{Exercise 2.7.1c}  
\textit{(c) Consider the subsequences \( (s_{2n}) \) and \( (s_{2n+1}) \), and show how the Monotone Convergence Theorem leads to a third proof for the Alternating Series Test.}
\medskip

\noindent\textbf{Proof.}
Let \((a_n)\) be a decreasing sequence of positive real numbers with \(a_n \to 0\).
Consider the two subsequences:
\[
s_{2n} = \sum_{k=1}^{2n} (-1)^{k+1} a_k = a_1 - a_2 + a_3 - \cdots - a_{2n}
\]
\[
s_{2n+1} = \sum_{k=1}^{2n+1} (-1)^{k+1} a_k = a_1 - a_2 + a_3 - \cdots + a_{2n+1}
\]

Note that \( s_{2(n+1)} - s_{2n} = a_{2n+1} - a_{2n+2} \ge 0 \), since \(a_n\) is decreasing. Hence \( (s_{2n}) \) is increasing.

Similarly, \( s_{2n+1} - s_{2(n+1)+1} = a_{2n+2} - a_{2n+3} \ge 0 \), so \( (s_{2n+1}) \) is decreasing.

Also, \( s_{2n+1} - s_{2n} = a_{2n+1} \ge 0 \), so \( s_{2n} \le s_{2n+1} \). Since \( a_n \to 0 \), we also have \( a_{2n+1} \to 0 \), so this difference goes to 0.

Since \( (s_{2n}) \) is increasing, and since \( s_{2n} \le s_{2n+1} \le s_1 \), it is bounded above. Similarly, \( (s_{2n+1}) \) is decreasing and bounded below by \( s_2 \).

By the Monotone Convergence Theorem, both \( (s_{2n}) \) and \( (s_{2n+1}) \) converge.

Let \( \lim s_{2n} = L \) and \( \lim s_{2n+1} = M \). Since \( s_{2n+1} - s_{2n} = a_{2n+1} \to 0 \), it follows that \( M - L = 0 \), so \( L = M \).

Because both the even and odd subsequences converge to the same limit \( L \), the full sequence \( (s_n) \) converges to \( L \).

\(\therefore \sum_{n=1}^{\infty} (-1)^{n+1} a_n\) converges.

\hfill$\Box$

\newpage

\noindent\textbf{Exercise 3.3.9}  
\textit{(c) Because \( x \in K \), there must exist an open set \( O_{\lambda_0} \) from the original collection that contains \( x \) as an element. Explain how this leads to the desired contradiction.}
\medskip

\noindent\textbf{Proof.}

Suppose, for contradiction, that \(K\) is not compact.

Since \(x \in K\) and \(\{O_\lambda\}_{\lambda \in \Lambda}\) covers \(K\), there is some \(\lambda_0\) such that \(x \in O_{\lambda_0}\).

Because \(O_{\lambda_0}\) is open and contains \(x\), there exists \(\varepsilon > 0\) such that \((x - \varepsilon, x + \varepsilon) \subseteq O_{\lambda_0}\).

As \(|I_n| \to 0\) by part (a), we can choose \(N\) so that for all \(n \ge N\), \(|I_n| < \varepsilon\).

If we write \(I_n = [a_n, b_n]\), then since \(x \in I_n\) and \(b_n - a_n = |I_n| < \varepsilon\), we have
\[
a_n > x - \varepsilon \quad \text{and} \quad b_n < x + \varepsilon.
\]

Then \(I_n \subseteq (x - \varepsilon, x + \varepsilon)\) for all \(n \ge N\).

Thus \(I_n \subseteq O_{\lambda_0}\) for all \(n \ge N\), so \(I_n \cap K \subseteq O_{\lambda_0}\). In particular, \(I_n \cap K\) is covered by just one open set.

This contradicts part (a), which says \(K \cap I_n\) has no finite subcover from \(\{O_\lambda\}_{\lambda \in \Lambda}\).

Therefore, our initial assumption that \(K\) is not compact leads to a contradiction. Hence every open cover of \(K\) has a finite subcover, and \(K\) is compact.

$\therefore$ \(K\) is compact.

\hfill$\Box$

\newpage

\noindent\textbf{Exercise 4.2.3}  
\textit{(c) Conjecture and verify $\lim_{x \to 1} t(x)$}
\medskip

\noindent\textbf{Proof.}

We conjecture
\[
\lim_{x \to 1} t(x) = 0.
\]

Let \(\varepsilon > 0\). Suppose \(t(x) \geq \varepsilon\). Writing \(x = \frac{m}{n}\) in lowest terms (with \(n \in \mathbb{N}\)), we have
\[
t(x) = \frac{1}{n} \geq \varepsilon \quad \Rightarrow \quad n \leq \frac{1}{\varepsilon}.
\]

Thus, if \(t(x) \geq \varepsilon\), the denominator \(n\) is bounded above by \(\frac{1}{\varepsilon}\).  
Consider only those \(x\) in the interval \((0,2)\). For each such \(n\), there are only finitely many rational numbers \(\frac{m}{n}\) in \((0,2)\) with this denominator. Hence the set
\[
E_\varepsilon = \{x \in (0,2) : t(x) \geq \varepsilon\}
\]
is finite. (Note: \(1 \in E_\varepsilon\) because \(t(1)=1 \geq \varepsilon\) for any \(\varepsilon \leq 1\); for \(\varepsilon >1\) the set is empty.)

If \(\varepsilon > 1\), then \(E_\varepsilon\) is empty, and the conclusion follows trivially. So we may assume \(0 < \varepsilon \le 1\); in this case \(E_\varepsilon\) is finite and contains \(1\).

Since \(E_\varepsilon\) is finite, let
\[
\delta = \min\{1,\; |1 - y| : y \in E_\varepsilon,\ y \neq 1\}.
\]
If \(E_\varepsilon = \{1\}\), simply take \(\delta = 1\). Then \(\delta > 0\) and by construction, if \(0 < |x-1| < \delta\), we have \(x \notin E_\varepsilon\) and \(x \in (0,2)\).

For any such \(x\):
\begin{itemize}
    \item If \(x\) is irrational, then \(t(x)=0<\varepsilon\).
    \item If \(x\) is rational, write \(x = \frac{m}{n}\) in lowest terms. Then \(x \in (0,2)\) and \(x \notin E_\varepsilon\) implies \(t(x) < \varepsilon\).
\end{itemize}
Thus \(t(x) < \varepsilon\) whenever \(0 < |x-1| < \delta\).

$\therefore \lim_{x \to 1} t(x) = 0.$

\hfill$\Box$

\newpage

\noindent\textbf{Exercise 4.5.3.}  
\textit{A function $f$ is increasing on $A$ if $f(x) \leq f(y)$ for all $x < y$ in $A$. Show that if $f$ is increasing on $[a,b]$ and satisfies the Intermediate Value Property (Definition 4.5.3), then $f$ is continuous on $[a,b]$.} 
\medskip

\noindent\textbf{Proof.}

Let \(c \in [a,b]\). We show \(f\) is continuous at \(c\).

\medskip
\noindent\textbf{Case 1: \(f(a) < f(c) < f(b)\).}

Let \(\varepsilon > 0\). Without loss of generality, assume
\[
\varepsilon < \min\{f(c) - f(a), f(b) - f(c)\}.
\]
Since any $\varepsilon>0$ can be reduced to satisfy this bound, proving continuity for such $\varepsilon$ is sufficient. Then,
\[
f(c) - \frac{\varepsilon}{2} > f(a) \quad \text{and} \quad f(c) + \frac{\varepsilon}{2} < f(b).
\]

Since \(f\) satisfies the Intermediate Value Property, there exist
\[
c_1 \in (a,c) \quad \text{with} \quad f(c_1) = f(c) - \frac{\varepsilon}{2}
\]
and
\[
c_2 \in (c,b) \quad \text{with} \quad f(c_2) = f(c) + \frac{\varepsilon}{2}.
\]

Because \(f\) is increasing, for any \(x\) with \(c_1 < x < c_2\) we have
\[
f(c_1) \le f(x) \le f(c_2).
\]
Thus
\[
f(c) - \frac{\varepsilon}{2} \le f(x) \le f(c) + \frac{\varepsilon}{2},
\]
so \(|f(x) - f(c)| \le \frac{\varepsilon}{2} < \varepsilon\).

Now set \(\delta = \min\{c - c_1, c_2 - c\}\). If \(|x - c| < \delta\), then \(c_1 < x < c_2\), and therefore \(|f(x) - f(c)| < \varepsilon\). Hence \(f\) is continuous at \(c\).

\medskip
\noindent\textbf{Case 2: \(f(c) = f(a)\).}

Let \(\varepsilon > 0\).

If \(f(b) = f(c)\), monotonicity implies \(f\) is constant on \([c,b]\), so \(|f(x)-f(c)| = 0\) for \(x \ge c\). For \(x \le c\), we will show below that \(f(x) = f(c)\). Thus \(f\) is continuous at \(c\) in this subcase. So assume \(f(b) > f(c)\).

Let \(\varepsilon' = \min\{\varepsilon,\; 2(f(b)-f(c))\}\). Then \(\varepsilon'>0\) and \(f(c)+\varepsilon'/2 \le f(b)\). By the Intermediate Value Property on \([c,b]\), there exists \(c_2 \in (c,b)\) with
\[
f(c_2) = f(c) + \frac{\varepsilon'}{2}.
\]

For any \(x \in (c, c_2)\), monotonicity gives \(f(c) \le f(x) \le f(c_2)\), so
\[
|f(x) - f(c)| \le \frac{\varepsilon'}{2} < \varepsilon'.
\]

For \(x \in [a,c)\), monotonicity gives \(f(x) \le f(c)\). Since \(f(c) = f(a)\), we actually have
\[
f(a) \le f(x) \le f(c) = f(a),
\]
which makes \(f(x) = f(a) = f(c)\). Hence \(f\) is constant on \([a,c]\), and \(|f(x)-f(c)| = 0 < \varepsilon\) for all \(x \in [a,c]\).

Set \(\delta = c_2 - c\). If \(|x - c| < \delta\) and \(x \ge c\), then \(x \in (c, c_2)\) and \(|f(x)-f(c)| < \varepsilon\). If \(x \le c\), then \(f(x) = f(c)\), so \(|f(x)-f(c)| = 0 < \varepsilon\). Thus \(f\) is continuous at \(c\).

\medskip
\noindent\textbf{Case 3: \(f(c) = f(b)\).}

Let \(\varepsilon > 0\).

If \(f(a) = f(c)\), then \(f\) is constant on \([a,c]\) by monotonicity, so \(|f(x)-f(c)| = 0\) for \(x \le c\). For \(x \ge c\), we will show \(f(x) = f(c)\). Hence continuity holds in this subcase. So assume \(f(a) < f(c)\).

The argument is symmetric to Case 2. Let \(\varepsilon' = \min\{\varepsilon,\; 2(f(c)-f(a))\}\). Then \(\varepsilon'>0\) and \(f(c)-\varepsilon'/2 \ge f(a)\). By the Intermediate Value Property on \([a,c]\), there exists \(c_1 \in (a,c)\) with
\[
f(c_1) = f(c) - \frac{\varepsilon'}{2}.
\]

For any \(x \in (c_1, c)\), monotonicity gives \(f(c_1) \le f(x) \le f(c)\), so
\[
|f(x) - f(c)| \le \frac{\varepsilon'}{2} < \varepsilon'.
\]

For any \(x \in (c,b]\), since \(f\) is increasing, \(f(c) \le f(x) \le f(b) = f(c)\), so \(f(x) = f(c)\). Thus \(|f(x) - f(c)| = 0 < \varepsilon\) for all \(x \in [c,b]\).

Take \(\delta = c - c_1\). If \(|x - c| < \delta\) and \(x \le c\), then \(x \in (c_1, c)\) and \(|f(x)-f(c)| < \varepsilon\). If \(x \ge c\), then \(x \in [c,b]\) and \(f(x) = f(c)\), so \(|f(x)-f(c)| = 0 < \varepsilon\). Hence \(f\) is continuous at \(c\).

$\therefore$ Since \(c \in [a,b]\) was arbitrary, \(f\) is continuous on \([a,b]\).

\hfill$\Box$

\end{document}
